\documentclass[10pt]{article}

\usepackage[margin=0.9in]{geometry}
\usepackage[T1]{fontenc}
\usepackage[utf8]{inputenc}
\usepackage{lmodern}
\usepackage{microtype}
\usepackage{amsmath,amssymb,mathtools}
\usepackage{booktabs}
\usepackage{tabularx}
\usepackage{array}
\usepackage{xcolor}
\usepackage{enumitem}
\usepackage[hidelinks]{hyperref}
\usepackage[capitalize,noabbrev]{cleveref}
\usepackage{caption}
\usepackage{xspace}
\usepackage{placeins}
\usepackage{longtable}

\definecolor{observedblue}{HTML}{075FA8}
\definecolor{placeholderred}{HTML}{B42318}
\definecolor{goodgreen}{HTML}{277A4B}

% Same convention as the main draft: blue = measured in this project, red = not yet filled.
\newcommand{\obs}[1]{\textcolor{observedblue}{#1}}
\newcommand{\tbd}[1]{\textcolor{placeholderred}{\textbf{[TBD: #1]}}}
\newcommand{\method}{\textsf{CAFE}\xspace}
\newcommand{\R}{\mathbb{R}}
\newcommand{\E}{\mathbb{E}}

\setlist[itemize]{leftmargin=*,topsep=2pt,itemsep=1pt}
\setlist[enumerate]{leftmargin=*,topsep=2pt,itemsep=1pt}
\captionsetup{font=small,labelfont=bf}

\title{\textbf{Can a cheap signal tell which generated image is better?}\\
\large Verification probes on Stable Diffusion 3.5-Medium: methods and results}
\author{}
\date{}

\begin{document}
\maketitle
\vspace{-0.4in}

\begin{abstract}
\noindent
A text-to-image model asked for ``a red cube and a blue sphere'' will sometimes produce it and
sometimes swap the colours. If we sample four images, one is usually better than the rest. A
\emph{verifier} is anything that guesses which one, without being told. This report measures how well
several candidate verifiers do that, on a frozen Stable Diffusion 3.5-Medium test bed, against
official automatic evaluators as ground truth. Two families are tested. The first,
\method, is \emph{decoding-free}: it compares the generator's own flow-matching residual under the
requested prompt against a minimally edited counterfactual prompt. It recovers \obs{7.4\%} of the
available headroom and is beaten, at established significance, by two controls that are simpler and
cheaper --- including one that uses no counterfactual at all. The second family scores each candidate
by its similarity to the real photograph the caption was written for. There, similarity in a
\emph{visual encoder} space recovers \obs{42.6\%}, similarity in the raw \emph{VAE latent} recovers
\obs{18.1\%} with a within-prompt rank correlation indistinguishable from zero, and similarity in the
\emph{denoiser's own intermediate hidden state} recovers \obs{38.8\%} with no decoder at all. Analyses
that compare whole representation spaces --- debiased CKA and SVCCA --- explain why: the hidden state
contains a linearly recoverable subspace shared with DINOv2 that the raw latent does not, and that
subspace does not live in the high-variance directions, which is why truncating to 99.9\% of variance
destroys it. We then train a small read-out head on the frozen denoiser, reaching cosine
\obs{0.586} to DINOv2 patch tokens, and show it carries \obs{36.7\%} of headroom without a decoder,
which is \obs{83\%} of what the target space itself achieves. Finally we ask how much of that target
space is needed: the alignment signal concentrates \emph{faster} than variance --- \obs{90\%} of the
full-rank result is retained at rank \obs{64}, where only \obs{45\%} of the spectral mass has
accumulated --- but it does not saturate, because the DINOv2 spectrum on this corpus is heavy-tailed
(\obs{$\alpha=0.800$}). The report defines every quantity, states every control, and reports what
each probe is blind to.
\end{abstract}

\section{What question is being asked, and why it matters}
\label{sec:why}

\subsection{The three-stage decomposition}

Improving a fast generator by selection and retraining decomposes into three questions that are
easily conflated:
\[
  \underbrace{\text{does a better sample exist?}}_{\text{headroom}}
  \;\longrightarrow\;
  \underbrace{\text{can we find it?}}_{\text{verification}}
  \;\longrightarrow\;
  \underbrace{\text{can the model absorb it?}}_{\text{amortization}}
\]
Each stage can fail independently, and a single end-to-end number cannot tell you which one did.
This report is entirely about the \textbf{middle} stage. It takes the candidate sets as given and
asks only: given four images produced for the same prompt, how well does a proposed cheap signal
reproduce the ordering that an expensive evaluator would give?

\subsection{Why the correlation is the whole point}

Everything hinges on a substitution. The expensive signal --- a vision--language model such as
VQAScore, or an official benchmark evaluator --- requires decoding every candidate latent to pixels
and then running a large scorer over it. If we want to \emph{train} on preferences over many
candidates, that cost is paid on every sample of every step.

A cheap signal is only useful if its ordering agrees with the expensive one. That agreement is a
correlation, and it is the object of this report for three reasons.

\begin{enumerate}
  \item \textbf{It is the thing that transfers.} A label is used by selecting or weighting
    candidates. Only the \emph{ordering} matters, not the absolute value, so rank agreement is the
    quantity of interest rather than, say, mean squared error against the evaluator.
  \item \textbf{It converts directly into an achievable gain.} \Cref{sec:headroom} defines
    recovered headroom $\rho$, which maps a correlation onto ``what fraction of the available
    improvement does this signal actually capture''. A verifier with a real but weak correlation
    can be worth less than its cost.
  \item \textbf{A weak correlation localizes the failure.} If a cheap signal fails, the question is
    whether the information is absent from the representation or merely inaccessible to the
    particular readout. \Cref{sec:repmetrics} gives tools that separate these, and
    \Cref{sec:refresults} shows a case where the answer is ``present, but not where you would look''.
\end{enumerate}

\subsection{The test bed, for a reader who knows nothing about it}

Everything below uses one frozen generator, Stable Diffusion 3.5-Medium (SD3.5-M), a
2.24B-parameter rectified-flow transformer operating on a $16\times64\times64$ VAE latent for a
$512\times512$ image. Candidates are produced by the \emph{unmodified} model at 8 sampling steps with
classifier-free guidance 7.0, four random seeds per prompt ($K=4$). No training of the generator
occurs anywhere in this report; the only thing ever trained is a small read-out head
(\Cref{sec:projector}).

Two prompt sets are used, and they are not interchangeable:
\begin{itemize}
  \item \textbf{T2I-CompBench++} prompts (\Cref{sec:cafe}) are short compositional templates
    (``a white piano and a black bench''). They have \emph{no} ground-truth photograph, so only
    prompt-based evaluators apply.
  \item \textbf{COCO captions} (\Cref{sec:refsim}) describe a real photograph, which is on disk.
    That makes \emph{reference-based} verifiers testable, which is impossible on CompBench++.
\end{itemize}

\section{Headroom and recovery: the measuring stick}
\label{sec:headroom}

\subsection{Definitions}

Fix a prompt $p$ with $K$ cached candidate latents $z_1,\dots,z_K$ and a pre-specified evaluator
$R(\cdot)$ applied to the decoded image. \emph{Headroom} is the gain available from perfect
selection over uniform selection:
\begin{equation}
  \mathrm{headroom}(p) \;=\; \max_i R(z_i) \;-\; \frac{1}{K}\sum_{i=1}^{K} R(z_i),
  \qquad
  H \;=\; \E_p\!\left[\mathrm{headroom}(p)\right].
\label{eq:headroom}
\end{equation}
The baseline is the \emph{mean} of the candidates, not the worst and not a single draw. This is
deliberate: a verifier that picks uniformly at random scores zero in expectation, because
$\E[R(\text{random pick})] = \frac1K\sum_i R(z_i)$.

A verifier assigns scores $S_i$ and selects $\hat\imath = \arg\max_i S_i$. Its quality is the
fraction of headroom it recovers:
\begin{equation}
  \rho \;=\;
  \frac{\sum_p \left[\, R(z_{\hat\imath(p)}) - \frac1K\sum_i R(z_i(p)) \,\right]}
       {\sum_p \left[\, \max_i R(z_i(p)) - \frac1K\sum_i R(z_i(p)) \,\right]}.
\label{eq:rho}
\end{equation}

\subsection{Four properties worth stating explicitly}

\begin{description}[leftmargin=1.2em,style=nextline]
  \item[It is a ratio of sums, not a mean of ratios.]
    Each prompt is weighted by its own headroom, so prompts where selection actually matters
    dominate. The alternative --- averaging the per-prompt fraction --- weights every prompt equally
    and is undefined when a prompt has zero headroom. The two differ materially: on the same data
    and the same verifier, ratio-of-sums gives \obs{42.59\%} where mean-of-ratios gives
    \obs{23.91\%}. Every number in this report is ratio-of-sums.
  \item[Prompts where all $K$ candidates tie are free.]
    They contribute $0$ to the numerator and $0$ to the denominator, so excluding them cannot change
    $\rho$. This is why such rows are skipped to save compute, and why the count is reported
    separately rather than hidden.
  \item[The oracle is exactly $1$ and chance is exactly $0$.] By construction.
  \item[It is \emph{not} antisymmetric.]
    $\rho$ of the reversed score is not $-\rho$, because for $K>2$ picking the minimum is not the
    mirror image of picking the maximum. This is why a verifier at $\obs{+7.38\%}$ has an inverted
    twin at $\obs{-10.73\%}$ rather than $-7.38\%$.
\end{description}

\subsection{The denominator differs between experiments --- a real trap}

$R$ is whatever evaluator that experiment froze, so $H$ is not one number.

\begin{table}[h]
\centering
\caption{Headroom denominators actually in use. Recoveries computed against different rows of this
table are \emph{not} comparable, because both the evaluator and the prompt pool differ.}
\label{tab:denominators}
\small
\begin{tabularx}{\textwidth}{@{}Xllr@{}}
\toprule
Experiment & Evaluator $R$ & Prompt pool & $H$ \\
\midrule
\method probe (\Cref{sec:cafe}) & per-category T2I-CompBench++ & 2,098 held-out & \obs{0.14067} \\
Reference similarity (\Cref{sec:refsim}) & VQAScore & 500 COCO captions & \obs{0.08753} \\
Historical ``best-of-4'' headline & CompBench++, \emph{selected by VQAScore} & same 2,098 & \obs{0.05852} \\
\bottomrule
\end{tabularx}
\end{table}

The third row deserves a warning. It selects the candidate with VQAScore and then \emph{measures} it
with CompBench++, so it is the gain of a particular cross-evaluator policy, not the CompBench++ oracle
headroom. Selecting with CompBench++ itself on identical rows gives \obs{0.14067}, which is
$2.4\times$ larger. Both quantities are legitimate; they answer different questions.

\section{Experiment 1: a decoding-free counterfactual verifier (\method)}
\label{sec:cafe}

\subsection{The idea, and the intuition for why it might work}

The generator is a conditional velocity field $v_\phi(x,\tau,p)$ trained so that, for a clean latent
$z$, Gaussian noise $\epsilon$, and interpolation level $\sigma$,
\[
  x_\sigma = (1-\sigma)z + \sigma\epsilon,
  \qquad
  v_\phi\big(x_\sigma, \tau(\sigma), p\big) \;\approx\; \epsilon - z .
\]
Define the residual
\begin{equation}
  \ell_\phi(z,p;\sigma,\epsilon)
  = \tfrac1d\big\| v_\phi(x_\sigma,\tau(\sigma),p) - (\epsilon-z) \big\|_2^2 .
\label{eq:residual}
\end{equation}
The intuition: if $z$ really depicts what $p$ describes, the model should find $z$ easy to predict
under $p$. The obvious problem is that $\ell$ also depends on how intrinsically hard $z$ is ---
a cluttered image is harder for any prompt. So the absolute residual conflates ``matches the
prompt'' with ``is a complicated picture''.

The fix is a \emph{contrast}. Build a counterfactual prompt $q$ that changes exactly one
compositional fact, and subtract:
\begin{equation}
  \delta_m = \ell_\phi(z,q;\sigma_m,\epsilon_m) - \ell_\phi(z,p;\sigma_m,\epsilon_m),
  \qquad
  B(z;p,q) = \frac{1}{M}\sum_{m=1}^{M}\delta_m .
\label{eq:contrast}
\end{equation}
Crucially both terms use the \textbf{same} $\sigma_m$ and the \textbf{same} noise draw
$\epsilon_m$ --- common random numbers. The image-difficulty term is then shared between the two
evaluations and cancels, leaving (in principle) only the effect of the changed fact. Positive $B$
favours $p$.

We use $M=8$ stratified points $\sigma_m = 0.05 + 0.90\,(m-\tfrac12)/M$ on $[0.05,0.95]$, avoiding
the degenerate endpoints, with $\tau = 1000\sigma$ passed to the transformer in float32.
Classifier-free guidance is off: CFG is a sampling heuristic, not the conditional field whose
residual we want.

\subsection{The counterfactual edit engine}

The quality of $q$ is the whole experiment. A sloppy edit measures vocabulary difficulty rather than
composition. Two families were built.

\begin{description}[leftmargin=1.2em,style=nextline]
  \item[C1, same-token-multiset swaps (primary).]
    The edit permutes words already present, so prompt length, token multiset and lexical difficulty
    are all held fixed and only the \emph{binding} changes. Three tiers:
    \begin{itemize}
      \item \texttt{c1\_attr}: swap two attribute values bound to \emph{different} head nouns.\\
        ``a red cube and a blue sphere'' $\rightarrow$ ``a blue cube and a red sphere''.
      \item \texttt{c1\_role}: swap the two noun phrases across an \emph{asymmetric} relation.\\
        ``a chicken behind a person'' $\rightarrow$ ``a person behind a chicken''.
      \item \texttt{c1\_count}: swap two different \emph{plural} counts, preserving number agreement.
    \end{itemize}
  \item[C2, single-token lexical replacement (secondary).]
    The pre-existing engine: replace one word with a different word (``blue'' $\rightarrow$
    ``purple''). This changes the vocabulary, so it cannot separate binding from lexical difficulty.
\end{description}

Several guards were necessary, and each was added because its absence produced a wrong edit:
symmetric relations (``next to'', ``near'') are refused because swapping across them does not change
the meaning; attribute pairs binding to the \emph{same} head noun are refused (``a fabric blanket and
a fluffy blanket'' describes the same scene either way); noun phrases containing a verb are refused
because the phrase boundary is not where a regular expression puts it; and article agreement is
repaired only where the swap \emph{created} the error, so that a counterfactual is never made more
fluent than an already-ungrammatical source prompt.

\begin{table}[h]
\centering
\caption{C1 edit-engine coverage over the 2,098 scored prompts. Headroom coverage weights each
prompt by its own oracle headroom, which is the quantity that matters for an unconditional
recovery.}
\label{tab:c1coverage}
\small
\begin{tabular}{@{}lrrrl@{}}
\toprule
Category & Rows covered & Row coverage & Headroom coverage & Tier \\
\midrule
texture     & \obs{296/300} & \obs{98.7\%} & \obs{98.3\%} & \texttt{c1\_attr} \\
3d\_spatial & \obs{296/300} & \obs{98.7\%} & \obs{99.2\%} & \texttt{c1\_role} \\
color       & \obs{272/300} & \obs{90.7\%} & \obs{85.3\%} & \texttt{c1\_attr} \\
shape       & \obs{268/300} & \obs{89.3\%} & \obs{91.0\%} & \texttt{c1\_attr} \\
spatial     & \obs{213/298} & \obs{71.5\%} & \obs{75.9\%} & \texttt{c1\_role} \\
complex     & \obs{147/300} & \obs{49.0\%} & \obs{54.3\%} & \texttt{c1\_attr} \\
numeracy    & \obs{133/300} & \obs{44.3\%} & \obs{41.6\%} & \texttt{c1\_count} \\
\midrule
\textbf{total} & \obs{\textbf{1,625/2,098}} & \obs{\textbf{77.5\%}} & \obs{\textbf{78.6\%}} & \\
\bottomrule
\end{tabular}
\end{table}

Coverage is low for numeracy because many prompts contain a single count (``one drum''), and for
spatial because symmetric relations are excluded by design. A defect in the pre-existing C2 engine
is recorded here because it affects earlier work: it selects the replacement as the next unused
entry of a flat family list, and those lists mix synonyms, so \obs{133 of 1,625} C2
``counterfactuals'' are meaning-preserving (\texttt{metallic}$\rightarrow$\texttt{metal},
\texttt{round}$\rightarrow$\texttt{circular}) --- \obs{25\%} of texture and \obs{20\%} of shape.

\subsection{Arms: every control, and what each one rules out}

All arms rank the same four latents with the same noise, so all comparisons are paired.

\begin{table}[h]
\centering
\caption{The arms and the hypothesis each one kills.}
\label{tab:arms}
\small
\begin{tabularx}{\textwidth}{@{}lX@{}}
\toprule
Arm & What it tests \\
\midrule
\texttt{cafe\_c1} & the proposed statistic: paired contrast under a same-multiset swap \\
\texttt{cafe\_c2} & does the \emph{edit family} matter? \\
\texttt{cafe\_c1\_selfnorm} & does dividing by cross-timestep dispersion help? \\
\texttt{cafe\_c1\_inverted} & orientation control: is the sign meaningful? \\
\texttt{cafe\_uncontrolled} & contrast against a fixed unrelated prompt (``a photograph''). \textbf{If this
  wins, the careful counterfactual is doing no work.} \\
\texttt{raw\_residual\_high/low} & $\pm\ell_\phi(z,p)$ alone. \textbf{If this wins, no counterfactual
  is needed at all}, and the cost halves \\
\texttt{clip} & decoded CLIPScore, the matched paid-for baseline \\
\texttt{random}, \texttt{oracle} & the two ends of the scale \\
\bottomrule
\end{tabularx}
\end{table}

\subsection{Results}

\begin{table}[h]
\centering
\caption{\method probe, full C1 set, $n=\obs{1{,}492}$, $M=8$, \obs{1.75}\,s per record on one H200.
Ground truth is the per-category CompBench++ evaluator, $H=\obs{0.14067}$.}
\label{tab:caferesults}
\small
\begin{tabular}{@{}lcrrr@{}}
\toprule
Verifier & Decode? & $\rho$ & 95\% CI & Top-1 \\
\midrule
Oracle & yes & \obs{100\%} & --- & \obs{100\%} \\
Decoded CLIP-L/14 & yes & \obs{30.18\%} & \obs{[+25.01, +35.10]} & \obs{32.8\%} \\
\texttt{cafe\_uncontrolled} & no & \obs{20.01\%} & \obs{[+14.40, +25.32]} & \obs{30.7\%} \\
\texttt{raw\_residual\_high} & no & \obs{16.22\%} & \obs{[+10.92, +21.50]} & \obs{32.8\%} \\
\texttt{cafe\_c2} & no & \obs{8.47\%} & \obs{[+2.78, +14.12]} & \obs{27.4\%} \\
\texttt{cafe\_c1\_selfnorm} & no & \obs{7.72\%} & \obs{[+2.14, +13.17]} & \obs{27.3\%} \\
\textbf{\texttt{cafe\_c1}} & no & \obs{\textbf{7.38\%}} & \obs{[+2.03, +12.84]} & \obs{26.4\%} \\
\texttt{cafe\_c1\_inverted} & no & \obs{$-10.73\%$} & \obs{[$-16.53$, $-4.97$]} & \obs{22.5\%} \\
\texttt{raw\_residual\_low} & no & \obs{$-27.90\%$} & \obs{[$-34.06$, $-21.88$]} & \obs{18.5\%} \\
Uniform & --- & \obs{$-0.14\%$} & \obs{[$-5.81$, $+5.22$]} & --- \\
\bottomrule
\end{tabular}
\end{table}

Marginal intervals overlap heavily and must not be read as a comparison. Because every verifier
ranks the same candidates on the same prompts, the paired difference is the correct test.

\begin{table}[h]
\centering
\caption{Paired prompt bootstrap, stratified by category, 20{,}000 draws, $n=\obs{1{,}492}$.
$\Delta$ is \texttt{cafe\_c1} minus the named arm.}
\label{tab:cafepaired}
\small
\begin{tabular}{@{}lrrrl@{}}
\toprule
Contrast & $\Delta\rho$ & 95\% CI & $p$ & Verdict \\
\midrule
$-$ \texttt{clip} & \obs{$-22.80\%$} & \obs{[$-30.15$, $-15.52$]} & \obs{$<10^{-4}$} & \obs{\method worse} \\
$-$ \texttt{uncontrolled} & \obs{$-12.63\%$} & \obs{[$-19.40$, $-5.90$]} & \obs{0.0004} & \obs{\method worse} \\
$-$ \texttt{raw\_residual\_high} & \obs{$-8.85\%$} & \obs{[$-16.13$, $-1.71$]} & \obs{0.0145} & \obs{\method worse} \\
$-$ \texttt{cafe\_c2} & \obs{$-1.09\%$} & \obs{[$-7.75$, $+5.50$]} & \obs{0.758} & no separation \\
$-$ \texttt{selfnorm} & \obs{$-0.35\%$} & \obs{[$-6.75$, $+6.04$]} & \obs{0.914} & no separation \\
$-$ \texttt{raw\_residual\_low} & \obs{$+35.28\%$} & \obs{[$+26.50$, $+44.08$]} & \obs{$<10^{-4}$} & \obs{\method better} \\
\bottomrule
\end{tabular}
\end{table}

\paragraph{The pre-specified mechanism test fails.}
The controlled contrast is outperformed by the \emph{uncontrolled} negative and by the raw residual
with \emph{no counterfactual at all}. It is also not distinguishable from the C2 edit family, so the
choice of C1 as primary has no measured support.

\paragraph{The aggregate hides a strong tier effect.}

\begin{table}[h]
\centering
\caption{By edit tier. The extensions are not weak; they are absent or negative.}
\label{tab:tiers}
\small
\begin{tabular}{@{}lrrrrr@{}}
\toprule
Tier & $n$ & \texttt{cafe\_c1} $\rho$ & 95\% CI & \texttt{clip} & \texttt{uncontrolled} \\
\midrule
\texttt{c1\_attr} (primary) & \obs{983} & \obs{15.62\%} & \obs{[+8.5, +22.5]} & \obs{41.36\%} & \obs{22.62\%} \\
\texttt{c1\_count}          & \obs{125} & \obs{$-1.28\%$} & \obs{[$-15.9$, $+13.4$]} & \obs{22.93\%} & \obs{13.74\%} \\
\texttt{c1\_role}           & \obs{384} & \obs{$-4.28\%$} & \obs{[$-14.5$, $+6.0$]} & \obs{13.19\%} & \obs{17.22\%} \\
\bottomrule
\end{tabular}
\end{table}

Restricted to \texttt{c1\_attr}, the paired contrasts soften to ``does not outperform'' rather than
``is outperformed'': versus \texttt{uncontrolled} \obs{$-7.01\%$}, $p=\obs{0.091}$; versus
\texttt{raw\_residual\_high} \obs{$-6.75\%$}, $p=\obs{0.145}$; versus \texttt{clip}
\obs{$-25.74\%$}, $p<\obs{10^{-4}}$. Both statements are true of different row sets and both should
be quoted.

\paragraph{Why the contrast underperforms its own controls: a coherent mechanism.}
Order the arms by how similar the subtracted reference prompt is to $p$, and read the within-prompt
correlation with the evaluator:
\[
  \underbrace{\ell_\phi(z,p)\ \text{alone}}_{\text{no reference}}\ \obs{+0.157}
  \;>\;
  \underbrace{\text{unrelated negative}}_{\text{``a photograph''}}\ \obs{+0.147}
  \;>\;
  \underbrace{\text{same-multiset swap}}_{\text{maximally similar}}\ \obs{+0.108}.
\]
Monotone \emph{decreasing} in the similarity of the subtrahend. This is exactly what happens if
$\ell_\phi(z,q)$ is dominated by a candidate-specific difficulty term shared with $\ell_\phi(z,p)$:
the more minimal the edit, the more completely that term cancels --- and the more of what survives is
noise. \textbf{The intervention that makes the contrast clean is the same intervention that destroys
its signal-to-noise.} Consistent with this, $B$ is only \obs{0.35\%} of the residual magnitude, and
removing the raw-energy component from $B$ barely changes it (\obs{$+0.108 \to +0.091$}), so the
residual signal is small rather than merely contaminated.

\paragraph{A separate, well-powered positive result.}
Averaged over candidates, $B$ is reliably positive: mean \obs{$+8.67\times10^{-4}$}, positive for
\obs{79.8\%} of candidates. The frozen field \emph{does} prefer the true prompt over its
same-multiset swap. But only \obs{45.5\%} of $B$'s variance is within-prompt, and the within-prompt
rank correlation is \obs{$+0.060$}. In words: \emph{the model knows which prompt is correct on
average; it does not know whether this particular sample got the binding right.}

\paragraph{The control that decides what the positive result means.}
Because $p$ and $q$ differ only in word order, a positive $B$ is equally consistent with sensitivity
to \emph{binding} and sensitivity to \emph{word order}. The discriminating control is a
meaning-preserving reorder $r$ with the identical token multiset:
\begin{center}
\begin{tabular}{@{}ll@{}}
$p$ & ``a black pen and a blue ink'' \\
$q$ & ``a blue pen and a black ink'' \quad (binding changed) \\
$r$ & ``a blue ink and a black pen'' \quad (binding \emph{identical}, order changed) \\
\end{tabular}
\end{center}

\begin{table}[h]
\centering
\caption{Reorder null, $n=\obs{785}$ captions, \obs{3{,}140} candidates.}
\label{tab:reorder}
\small
\begin{tabular}{@{}lrrr@{}}
\toprule
 & mean $\delta$ & fraction $>0$ & $\rho$ as a selector \\
\midrule
counterfactual $q$ & \obs{$+8.67\times10^{-4}$} & \obs{79.8\%} & \obs{15.91\%} \\
reorder $r$ (meaning preserved) & \obs{$+2.38\times10^{-4}$} & \obs{69.5\%} & \obs{$-0.98\%$} \\
paired difference & \obs{$+6.29\times10^{-4}$} & \multicolumn{2}{l}{95\% CI \obs{$[5.83, 6.78]\times10^{-4}$}} \\
\bottomrule
\end{tabular}
\end{table}

The directional effect is therefore \obs{$\approx73\%$} binding and \obs{$\approx27\%$} word-order
sensitivity, and as a \emph{selector} the reorder is a proper null (\obs{$-0.98\%$}, versus
\obs{$+0.39\%$} for random). The binding component survives its control. Per category the binding
share is cleanest for colour (reorder only \obs{12.5\%} of the counterfactual effect), then texture
\obs{34.7\%}, and weakest for shape \obs{68.5\%}.

\section{Experiment 2: does resemblance to the real photograph predict prompt satisfaction?}
\label{sec:refsim}

\subsection{The idea, and why it is worth asking}

A COCO caption was written to describe a specific photograph, and that photograph is on disk. So for
these prompts, and only these, a \emph{reference-based} verifier is possible: score each candidate by
how much it resembles the ground truth.

This matters because at \emph{training} time a reference usually exists. If resemblance predicts
prompt satisfaction, a training-time label can be produced with no vision--language model and, in the
latent and hidden-state variants, no decoder.

\paragraph{The confound, stated first.}
One caption describes many valid photographs. Cosine to \emph{one} chosen reference measures
instance-level similarity --- layout, palette, viewpoint --- not prompt adherence. A weak correlation
is therefore the \emph{expected} result and would not show that the embedding lacks semantics. A
strong one is the surprising outcome.

\subsection{Setup}

\obs{500} COCO captions, each with its real photograph, and \obs{4} candidates each generated by the
frozen SD3.5-M at 8 steps and CFG 7.0 --- the same recipe as \Cref{sec:cafe}. Each candidate is
scored by VQAScore (\texttt{clip-flant5-xxl}) against its caption; $H=\obs{0.08753}$, with
\obs{zero} tied prompts. Captions whose \emph{text} repeats across different photographs were dropped
(\obs{12} of them), because otherwise a text join attaches a generation to the wrong ground truth.

\subsection{The four spaces, and why each is the natural thing to try}

\begin{description}[leftmargin=1.2em,style=nextline]
  \item[\texttt{dino} --- DINOv2-base CLS, 768-d, requires a decode.]
    A self-supervised visual encoder whose features are known to capture object identity and layout.
    The natural ``semantic similarity of two images'' measure.
  \item[\texttt{clip} --- CLIP ViT-L/14 image embedding, 768-d, requires a decode.]
    Trained on image--text pairs, so it is tilted towards exactly the caption-relevant content.
  \item[\texttt{latent} --- the SD3.5 VAE latent, $16\times64\times64 = 65{,}536$-d, decoding-free.]
    The cheapest possible option: the representation the generator already produces. Its weakness is
    that a VAE is trained for \emph{reconstruction}, so its latent is organised by
    appearance, not meaning.
  \item[\texttt{repa} --- the denoiser's hidden state at block 8, 1536-d, decoding-free.]
    The premise of representation-alignment methods (REPA) is that a diffusion transformer's
    intermediate activations, while denoising, align with a visual encoder's features. If true, this
    is a decoding-free representation with encoder-like semantics. It is extracted from the
    \emph{image} stream of the joint block --- the text stream is a different tensor and using it
    would silently measure the wrong thing --- while the model processes the image noised to
    $\sigma=0.5$, with reference and candidates sharing the same noise and the same caption
    conditioning so that only the image differs.
\end{description}

Each space is also evaluated inside a PCA basis fitted on that run's pooled reference-and-candidate
vectors and truncated at 99.9\% of variance. This is a rotation and truncation \emph{inside one
space}; nothing is fitted against the evaluator.

\paragraph{A note on what cannot be done.}
There is no cosine between a 768-d DINOv2 embedding and a 65{,}536-d latent. They have no shared
basis, and PCA cannot supply one: projecting means a dot product with a component vector, so the
mismatch sits on the \emph{input} side, where truncating the output dimension has no effect. The
tools in \Cref{sec:repmetrics} are what answer cross-space questions.

\subsection{Results}
\label{sec:refresults}

\begin{table}[h]
\centering
\caption{Comprehensive correlation table. $n=\obs{500}$ captions, $K=4$, ground truth VQAScore,
$H=\obs{0.08753}$. ``Within-prompt'' correlates the $K$ similarities against the $K$ evaluator
scores inside each prompt and averages over prompts --- this is the decisive form, because it asks
whether the signal picks the better \emph{candidate}. ``Pooled'' correlates all $n\!\times\!K$ pairs
at once and is confounded by prompt difficulty. Decoding-free arms in bold.}
\label{tab:refsim}
\small
\begin{tabular}{@{}lcrrrrrrr@{}}
\toprule
& & & & \multicolumn{2}{c}{Within-prompt} & & \multicolumn{2}{c}{Pooled} \\
\cmidrule(lr){5-6}\cmidrule(lr){8-9}
Arm & Dec.? & $\rho$ & 95\% CI & Spearman & 95\% CI & Pearson & Spearman & Pearson \\
\midrule
\texttt{dino\_patchmean} & yes & \obs{43.92\%} & \obs{[+33.7,+53.5]} & \obs{+0.2149} & \obs{[+.164,+.265]} & \obs{+0.2414} & \obs{+0.2099} & \obs{+0.3351} \\
\texttt{clip\_pca} & yes & \obs{43.64\%} & \obs{[+32.9,+53.4]} & \obs{+0.2435} & \obs{[+.194,+.291]} & \obs{+0.2864} & \obs{+0.2388} & \obs{+0.3172} \\
\texttt{dino}      & yes & \obs{42.59\%} & \obs{[+32.2,+52.2]} & \obs{+0.1852} & \obs{[+.135,+.233]} & \obs{+0.2343} & \obs{+0.1965} & \obs{+0.2575} \\
\texttt{clip}      & yes & \obs{41.89\%} & \obs{[+31.2,+52.0]} & \obs{+0.2554} & \obs{[+.206,+.305]} & \obs{+0.2744} & \obs{+0.1671} & \obs{+0.1977} \\
\texttt{dino\_pca} & yes & \obs{41.65\%} & \obs{[+31.2,+51.4]} & \obs{+0.1924} & \obs{[+.143,+.241]} & \obs{+0.2370} & \obs{+0.1975} & \obs{+0.2582} \\
\textbf{\texttt{repa}} & \textbf{no} & \obs{\textbf{38.75\%}} & \obs{[+28.3,+48.7]} & \obs{+0.1679} & \obs{[+.118,+.218]} & \obs{+0.1926} & \obs{+0.2173} & \obs{+0.3088} \\
\textbf{\texttt{repa\_projected}} & \textbf{no} & \obs{\textbf{36.66\%}} & \obs{[+25.6,+47.1]} & \obs{+0.1504} & \obs{[+.099,+.201]} & \obs{+0.2025} & \obs{+0.1602} & \obs{+0.3011} \\
\textbf{\texttt{latent}} & \textbf{no} & \obs{18.11\%} & \obs{[+5.3,+30.5]} & \obs{+0.0286} & \obs{[$-$.025,+.081]} & \obs{+0.0462} & \obs{+0.0231} & \obs{+0.0777} \\
\textbf{\texttt{repa\_pca}} & \textbf{no} & \obs{17.11\%} & \obs{[+4.5,+29.5]} & \obs{+0.0534} & \obs{[+.003,+.104]} & \obs{+0.0863} & \obs{+0.0694} & \obs{+0.0921} \\
\textbf{\texttt{latent\_pca}} & \textbf{no} & \obs{6.48\%} & \obs{[$-$7.0,+19.7]} & \obs{$-0.0219$} & \obs{[$-$.073,+.030]} & \obs{$-0.0213$} & \obs{+0.0049} & \obs{+0.0504} \\
\bottomrule
\end{tabular}
\end{table}

Three readings.

\begin{enumerate}
  \item \textbf{Resemblance to the ground truth genuinely predicts prompt satisfaction.} Despite the
    one-caption-many-images confound, encoder-space similarity recovers \obs{$\approx$43\%} of
    headroom. The plain interpretation is that a generation which fails the caption also tends to
    look unlike the photograph the caption describes.
  \item \textbf{The raw latent is close to inert.} Its within-prompt Spearman is \obs{$+0.029$} with
    a confidence interval straddling zero. This reproduces, in a reference-based setting, the
    project's recurring finding that latent-to-semantics readouts do not work.
  \item \textbf{The denoiser's hidden state is not inert, and needs no decoder.} \obs{38.75\%},
    within confidence of DINO and CLIP, both of which must decode. Same model, same images, one
    layer deeper: \obs{$+0.029 \to +0.168$}.
\end{enumerate}

\paragraph{PCA at 99.9\% variance destroys precisely the two decoding-free arms.}
\texttt{repa} \obs{$38.75\% \to 17.11\%$} and \texttt{latent} \obs{$18.11\% \to 6.48\%$}, while the
decoded arms are untouched. Retaining 99.9\% of variance still required \obs{2{,}376} latent
components and \obs{727} hidden-state components, so this is not aggressive truncation. The
conclusion is that in those two spaces the semantic content sits in \emph{low-variance} directions,
and variance-ranked truncation deletes it. This is a direct caution against the natural idea of
reducing everything to a fixed-size PCA embedding.

\section{Comparing whole representation spaces}
\label{sec:repmetrics}

The remaining question --- \emph{why} does the hidden state work where the latent does not --- is
about whole spaces of different dimensionality, so cosine cannot answer it. Three tools apply.
Implementations and self-tests are in \texttt{phaseN/repsim\_metrics.py}.

\subsection{Linear CKA, and the bias that makes the naive version unusable}

Centered Kernel Alignment compares the \emph{Gram matrices} of two representations:
\begin{equation}
  \mathrm{CKA}(X,Y) = \frac{\|Y^\top X\|_F^2}{\|X^\top X\|_F \, \|Y^\top Y\|_F},
\end{equation}
with $X\in\R^{n\times p}$, $Y\in\R^{n\times q}$ column-centred. Intuitively it asks: \emph{do the two
spaces consider the same pairs of samples to be similar?} Because it works through
sample-by-sample similarity, it never needs $p=q$. It is invariant to rotation and to isotropic
scaling, which is what makes it a fair comparison between arbitrary representations.

Its weakness is a positive bias when features outnumber samples. Verified on synthetic data:
two \emph{independent} Gaussians of representative shapes score \obs{$+0.179$}, not $0$. The unbiased
HSIC estimator removes this: the same independent pair scores \obs{$-0.004$}, while identical
representations still score $1$ and a planted shared signal scores \obs{$+0.870$}. The debiased value
is the one to read. The two algebraic forms of CKA (feature-space and Gram/HSIC) were checked to
agree to \obs{$2.2\times10^{-16}$}.

\paragraph{What CKA is blind to.} It is dominated by the high-variance directions, because those
dominate the Gram matrix. Two spaces that share their dominant structure will score high even if
they disagree about everything else.

\subsection{SVCCA, and the split that makes it honest}

SVCCA truncates each space by SVD at 99\% of variance and then runs canonical correlation analysis
between the reduced spaces, reporting the mean canonical correlation. Intuitively it asks a different
question: \emph{is there a linear subspace of $X$ that is a linear function of a subspace of $Y$?}
Unlike CKA, it does not care whether that subspace carries much variance --- only whether it is
linearly decodable.

CCA fitted and evaluated on the same rows finds large correlations between pure noise as soon as the
retained dimensions approach the sample count: measured at \obs{$r=0.296$} on independent Gaussians.
Directions are therefore fitted on one half of the samples and correlations measured on the held-out
half, with a row-permutation null reported beside every value. With that split, independent data
gives \obs{$0.054$} against a null of \obs{$0.061$} (excess \obs{$-0.007$}), and a planted shared
signal gives \obs{$0.303$} against \obs{$0.053$} (excess \obs{$+0.250$}).

\subsection{RSA}

Representational Similarity Analysis is the cheapest and most assumption-free: for each prompt, take
the $K$ reference-to-candidate similarities computed in space $A$ and in space $B$, and correlate
those two vectors of length $K$. It fits nothing at all and answers exactly the question the
downstream use cares about: \emph{would these two spaces rank the candidates the same way?}

\subsection{Results, and the informative disagreement}

\begin{table}[h]
\centering
\caption{Cross-space structure on the pooled \obs{2{,}500} reference and candidate vectors.
Debiased CKA is the value to read. SVCCA excess is held-out score minus permutation null.
RSA is the mean within-prompt Spearman between similarity structures.}
\label{tab:crossspace}
\small
\begin{tabular}{@{}lrrrrrr@{}}
\toprule
Pair & CKA (biased) & \textbf{CKA (debiased)} & SVCCA & null & \textbf{excess} & RSA \\
\midrule
\texttt{latent}--\texttt{repa} & \obs{+0.9211} & \obs{\textbf{+0.9213}} & \obs{+0.1534} & \obs{+0.0226} & \obs{+0.1308} & \obs{+0.4036} \\
\texttt{clip}--\texttt{dino}   & \obs{+0.5918} & \obs{\textbf{+0.5751}} & \obs{+0.1810} & \obs{+0.0239} & \obs{+0.1571} & \obs{+0.2992} \\
\texttt{clip}--\texttt{repa}   & \obs{+0.2389} & \obs{\textbf{+0.2355}} & \obs{+0.2921} & \obs{+0.0226} & \obs{\textbf{+0.2695}} & \obs{+0.2452} \\
\texttt{clip}--\texttt{latent} & \obs{+0.1975} & \obs{\textbf{+0.1899}} & \obs{+0.0326} & \obs{+0.0227} & \obs{+0.0099} & \obs{+0.0884} \\
\texttt{dino}--\texttt{repa}   & \obs{+0.1923} & \obs{\textbf{+0.1876}} & \obs{+0.2939} & \obs{+0.0251} & \obs{\textbf{+0.2689}} & \obs{+0.2004} \\
\texttt{dino}--\texttt{latent} & \obs{+0.1792} & \obs{\textbf{+0.1692}} & \obs{+0.0257} & \obs{+0.0221} & \obs{\textbf{+0.0036}} & \obs{+0.0512} \\
\midrule
\multicolumn{7}{@{}l}{\emph{involving the trained read-out head of \Cref{sec:projector}}}\\
\texttt{dino\_patchmean}--\texttt{repa\_projected} & \obs{+0.8272} & \obs{\textbf{+0.8275}} & \obs{+0.5687} & \obs{+0.0229} & \obs{\textbf{+0.5458}} & \obs{+0.5932} \\
\texttt{dino\_patchmean}--\texttt{repa} & \obs{+0.3230} & \obs{\textbf{+0.3208}} & \obs{+0.4143} & \obs{+0.0215} & \obs{+0.3928} & \obs{+0.3620} \\
\texttt{dino\_patchmean}--\texttt{latent} & \obs{+0.2723} & \obs{\textbf{+0.2669}} & \obs{+0.0345} & \obs{+0.0228} & \obs{\textbf{+0.0117}} & \obs{+0.1612} \\
\bottomrule
\end{tabular}
\end{table}

CKA and SVCCA give \emph{opposite orderings}, and that is the result rather than a problem.

\begin{itemize}
  \item CKA says \texttt{latent} and \texttt{repa} are nearly the same space (\obs{0.921}) and that
    \texttt{repa}--\texttt{dino} is low (\obs{0.188}).
  \item SVCCA says \texttt{repa}--\texttt{dino} has the \emph{largest} shared decodable subspace
    (excess \obs{$+0.269$}) while \texttt{latent}--\texttt{dino} has essentially none
    (excess \obs{$+0.004$}).
\end{itemize}

Both are correct about different things, and together they explain \Cref{tab:refsim}. The latent and
the layer-8 hidden state share their \emph{dominant variance} structure --- both are dominated by
low-level image content --- which is what CKA measures, hence \obs{0.92}. But the hidden state
additionally contains a \emph{linearly recoverable} DINOv2-aligned subspace that the raw latent does
not, which is what SVCCA measures, hence \obs{$+0.269$} versus \obs{$+0.004$}.

That is the representation-alignment premise, confirmed on this model without training anything. It
also explains why PCA at 99.9\% variance collapsed the \texttt{repa} arm: the DINO-aligned subspace
is real and linearly accessible, but it is not among the high-variance directions, so
variance-ranked truncation removes it.

\section{A trained read-out head, and the cross-space arm it enables}
\label{sec:projector}

The analyses above deliberately fit nothing. They therefore cannot answer the practical question:
how much of the visual encoder's information can be \emph{recovered} from the denoiser when a map is
allowed? That question matters because a map is exactly what a training-time labeller would use.

\subsection{Training}

We train a per-token MLP $1536 \rightarrow 2048 \rightarrow 2048 \rightarrow 768$ that maps the
frozen denoiser's layer-8 hidden state to DINOv2-base patch tokens, maximising cosine, on COCO
\texttt{train2017}. Only this head trains: the denoiser and the encoder are frozen, so nothing here
can change the generator and the result is purely a read-out. Validation uses a disjoint slice of
\texttt{train2017}, leaving \texttt{val2017} --- which supplies the reference photographs of
\Cref{sec:refsim} --- entirely untouched by training or early stopping.

Patch-to-patch alignment requires the two token grids to correspond, verified numerically before any
compute: at 512\,px the model emits $(512/8/2)^2 = 1024$ image tokens and DINOv2 at
$224\cdot\lfloor 512/256\rfloor = 448$\,px emits $(448/14)^2 = 1024$ patches.

\begin{table}[h]
\centering
\caption{Read-out head training. A randomly initialised projector sits near $1-\cos = 1$.
The curve is still descending at the final step, so \obs{0.586} is a lower bound on what this map
can reach.}
\label{tab:projector}
\small
\begin{tabular}{@{}lrrrrrrr@{}}
\toprule
step & 500 & 1500 & 2500 & 3500 & 4500 & 5500 & \textbf{6000} \\
\midrule
val $1-\cos$ & \obs{0.5115} & \obs{0.4649} & \obs{0.4459} & \obs{0.4399} & \obs{0.4269} & \obs{0.4209} & \obs{\textbf{0.4140}} \\
val $\cos$   & \obs{0.489} & \obs{0.535} & \obs{0.554} & \obs{0.560} & \obs{0.573} & \obs{0.579} & \obs{\textbf{0.586}} \\
\bottomrule
\end{tabular}
\end{table}

A frozen intermediate hidden state predicts DINOv2 patch tokens at cosine \obs{0.586}. This is the
direct confirmation of what \Cref{tab:crossspace} inferred from subspace statistics alone: the
encoder's information is present in the denoiser and a small head reaches it.

\subsection{The cross-space arm}

With a map in hand, the comparison that was previously undefined becomes well posed:
\[
  \texttt{repa\_projected} \;=\; \cos\!\Big(\operatorname{mean}_\text{tok} \text{proj}\big(h_T(\text{candidate})\big),\;
  \operatorname{mean}_\text{patch} \text{DINOv2}(\text{reference})\Big).
\]
Two choices decide whether this number means anything.

\paragraph{Pooling before the cosine is required, not a convenience.}
The head was trained per token against the patch grid of the \emph{same} image, so token $i$ of the
prediction and of the target are the same spatial location. A candidate and a reference are
different images and share no such correspondence, so a per-token cosine between them would compare
unrelated positions. Both sides are pooled over tokens first.

\paragraph{The matched upper bound must live in the target space.}
Comparing \texttt{repa\_projected} against the CLS-based \texttt{dino} arm would be unmatched,
because the head was trained to hit \emph{patch} tokens, not CLS. \Cref{tab:refsim} therefore also
reports \texttt{dino\_patchmean}, the same similarity computed between reference and candidate
patch-means. Without that row one could not tell whether a gap is the head's fault or the target
space's.

The result is \obs{36.66\%} against a matched decoded ceiling of \obs{43.92\%}: the learned map
carries \obs{83\%} of what the target space itself achieves, with no decoder. The structural
measures confirm the map did its job --- \texttt{dino\_patchmean}--\texttt{repa\_projected} reaches
debiased CKA \obs{0.828} and SVCCA excess \obs{+0.546}, against \obs{0.267} and \obs{+0.012} for
\texttt{dino\_patchmean}--\texttt{latent}.

Two honest qualifications. The head does \emph{not} beat the raw hidden state
(\obs{36.66\%} versus \obs{38.75\%}); mapping into the encoder's space costs a little. And the raw
hidden states are nearly collinear --- mean cosine \obs{0.970} --- so \texttt{repa}'s similarity is a
small perturbation on a large common component, whereas the projected arm at \obs{0.738} is better
conditioned even though its $\rho$ is slightly lower.

\section{How much of the visual target space is needed?}
\label{sec:rank}

A read-out head is cheaper, and an auxiliary alignment loss is better conditioned, when the target is
low-rank. The natural construction is to keep only the top-$n$ principal components of the encoder,
$z_n(x) = P_n\big(E_D(x) - \bar v\big)$. The question this raises is not whether a low-rank target
reconstructs the encoder --- by definition it captures the top-$n$ spectral mass --- but whether it
retains the \emph{alignment} signal. Those are different functionals of the same feature, and a
direction carrying little variance can still carry the binding information.

\paragraph{Protocol.}
The PCA basis is estimated once on \obs{5{,}000} COCO \texttt{train2017} images
(\obs{100{,}000} patch embeddings, dimension \obs{768}) and applied unchanged to \texttt{val2017}
references and to generated candidates, so nothing is fitted on the evaluation set. For each rank we
score every candidate by its cosine to the reference \emph{inside the rank-$n$ subspace} and
correlate with VQAScore. Patches are pooled before projection, which is exact rather than
approximate because the projection is linear (verified to \obs{$6.2\times10^{-15}$}).

\begin{table}[h]
\centering
\caption{Alignment recovery against retained spectral mass, $n=\obs{500}$ captions. Cumulative mass
is a property of the encoder alone; $\rho$ is measured against VQAScore.}
\label{tab:ranksweep}
\small
\begin{tabular}{@{}rrrrr@{}}
\toprule
rank $n$ & cumulative spectral mass & $\rho$ & 95\% CI & within-prompt Spearman \\
\midrule
4    & \obs{13.0\%} & \obs{27.43\%} & \obs{[+15.7, +38.9]} & \obs{+0.0860} \\
8    & \obs{18.9\%} & \obs{33.67\%} & \obs{[+22.3, +44.6]} & \obs{+0.1396} \\
16   & \obs{25.9\%} & \obs{34.38\%} & \obs{[+23.0, +45.0]} & \obs{+0.1641} \\
32   & \obs{34.4\%} & \obs{36.47\%} & \obs{[+25.1, +47.3]} & \obs{+0.1582} \\
48   & \obs{40.4\%} & \obs{38.81\%} & \obs{[+27.7, +48.9]} & \obs{+0.1733} \\
64   & \obs{45.2\%} & \obs{38.61\%} & \obs{[+27.6, +48.7]} & \obs{+0.1652} \\
128  & \obs{58.9\%} & \obs{42.62\%} & \obs{[+32.1, +52.2]} & \obs{+0.1694} \\
256  & \obs{75.3\%} & \obs{42.30\%} & \obs{[+32.2, +51.6]} & \obs{+0.1767} \\
512  & \obs{92.1\%} & \obs{39.19\%} & \obs{[+28.6, +48.9]} & \obs{+0.1736} \\
768  & \obs{100.0\%} & \obs{42.87\%} & \obs{[+32.8, +52.0]} & \obs{+0.1840} \\
\bottomrule
\end{tabular}
\end{table}

\paragraph{Alignment concentrates faster than variance.}
Rank \obs{64} retains \obs{90\%} of the full-rank recovery (\obs{38.61\%} of \obs{42.87\%}) while
only \obs{45.2\%} of the spectral mass has accumulated; rank \obs{16} already retains \obs{80\%} at
\obs{25.9\%} of mass. A low-rank target is therefore a cheap and largely lossless carrier of the
alignment signal, and this is a stronger justification for a small $n$ than a variance argument
would give.

\paragraph{But there is no elbow to saturate at.}
The DINOv2 spectrum on this corpus is heavy-tailed. Fitting $\lambda_i \approx C i^{-\alpha}$ over
eigenvalues 5--512 gives \obs{$\alpha = 0.800$}, and reaching \obs{90\%} of the mass requires rank
\obs{467}, \obs{99\%} rank \obs{718}. Correspondingly $\rho$ does not plateau at 64: it drifts
upward to rank $\sim$128 and the within-prompt rank correlation increases monotonically to full rank.
All confidence intervals overlap heavily, so rank 128 cannot be claimed to beat rank 64. The
defensible statement is the weaker and more useful one: \textbf{truncation is cheap in alignment
terms long before it is cheap in variance terms, but nothing saturates}, so a rank chosen by a
spectral-elbow argument would be chosen on a feature the spectrum does not have. Any claim that
performance saturates at a spectral elbow should be checked against $\alpha$ on the actual corpus,
because at $\alpha<1$ the tail mass does not decay at the rate such arguments assume.

\section{What this means for a decoding-free training signal}
\label{sec:implications}

The motivating goal is a label computable \emph{before} decoding, so that preference distillation
does not pay a VAE decode plus a large vision--language model per candidate per step.
\Cref{tab:status} summarises where that stands.

\begin{table}[h]
\centering
\caption{Status of each route to a pre-decoding label. $\rho$ values come from different pools and
evaluators (\Cref{tab:denominators}) and are not directly comparable across the horizontal rule.}
\label{tab:status}
\small
\begin{tabularx}{\textwidth}{@{}Xccr@{}}
\toprule
Route & Needs decode? & Needs a reference image? & $\rho$ \\
\midrule
Counterfactual flow-residual contrast (\Cref{sec:cafe}) & no & no & \obs{7.38\%} \\
Raw residual magnitude & no & no & \obs{16.22\%} \\
Uncontrolled negative prompt & no & no & \obs{20.01\%} \\
Decoded CLIPScore against the prompt & yes & no & \obs{30.18\%} \\
\midrule
VAE-latent similarity to reference (\Cref{sec:refsim}) & no & yes & \obs{18.11\%} \\
Read-out head into encoder space (\Cref{sec:projector}) & no & yes & \obs{36.66\%} \\
Hidden-state similarity to reference & no & yes & \obs{38.75\%} \\
Decoded encoder similarity to reference & yes & yes & \obs{43.92\%} \\
\bottomrule
\end{tabularx}
\end{table}

Three conclusions follow.

\begin{enumerate}
  \item \textbf{No reference-free pre-decoding label is available yet.} The best measured is
    \obs{20.01\%}, from an \emph{uncontrolled} negative prompt, and the mechanism analysis of
    \Cref{sec:cafe} indicates that much of it tracks image complexity rather than semantics: the raw
    residual alone correlates \obs{$+0.157$} with the evaluator, more than the carefully controlled
    contrast at \obs{$+0.108$}.
  \item \textbf{A reference-based pre-decoding label \emph{is} available and is competitive.} At
    \obs{38.75\%} it is statistically indistinguishable from decoded encoder similarity and above
    decoded CLIPScore's \obs{30.18\%} on its own pool. Since distillation on real-image corpora has a
    reference for every training prompt, this is directly usable as a training-time labeller.
  \item \textbf{The representation, not the modelling capacity, was the obstacle.} A latent-to-encoder
    read-out was previously judged inert. That judgement was reached on the \emph{VAE latent}, whose
    shared decodable subspace with the encoder is \obs{+0.004} --- indistinguishable from nothing. The
    denoiser's hidden state is a different object, at \obs{+0.269}, and a head trained on it reaches
    cosine \obs{0.586}. Conclusions drawn about the latent should not be carried over to the hidden
    state.
\end{enumerate}

\paragraph{The open question.} Every working route above needs a reference image, so none of them
labels a synthetic compositional prompt. The untested route that would is a probe trained directly
from the hidden state to the alignment score, with no reference and no decode. The labels for it
already exist: \obs{20{,}168} candidate-level scores on training prompts and \obs{8{,}392}
candidate-level official-evaluator scores on held-out prompts. Two design points matter. The metric
depends only on the ordering \emph{within} a prompt, so the objective should be a within-group
ranking loss rather than a regression onto absolute scores, which spends capacity on the
across-prompt difficulty that $\rho$ divides out. And the identical probe should be run on the raw
VAE latent as a matched control, since the analyses here predict it stays inert.

\tbd{direct hidden-state-to-score probe: held-out $\rho$ against the decoded-CLIP bar of
\obs{30.18\%}, with latent-probe and random-label controls.}

\section{Threats to validity}
\label{sec:threats}

\begin{itemize}
  \item \textbf{Round-trip latents.} In \Cref{sec:cafe} the latent $z$ is obtained by VAE-encoding a
    decoded PNG. It is a round trip, not the sampler's terminal latent.
  \item \textbf{Teacher, not student.} All candidates are 8-step CFG-7 samples from the frozen
    teacher, not from a fast distilled student.
  \item \textbf{Different denominators.} \Cref{tab:denominators} applies throughout; \obs{7.38\%} and
    \obs{42.59\%} are fractions of different quantities measured by different evaluators on different
    prompts, and must not be placed in one ranking.
  \item \textbf{Reference ambiguity.} \Cref{sec:refsim} scores against one photograph among many
    valid ones, so its correlations are a lower bound on what a reference-based signal could achieve
    with multiple references.
  \item \textbf{One backbone, one resolution.} Everything is SD3.5-M at $512^2$.
  \item \textbf{Layer and noise level not swept.} The \texttt{repa} arm uses block 8 at $\sigma=0.5$
    because that is the layer the project's alignment runs used. Neither choice was tuned, so the
    \obs{38.75\%} is not an optimum.
  \item \textbf{Abstention untested.} No arm in this report declines to answer, so the reported
    unconditional recovery is the conditional recovery times a constant coverage factor, by
    construction rather than by measurement.
\end{itemize}

\section{Reproduction}

\begin{table}[h]
\centering
\small
\begin{tabularx}{\textwidth}{@{}lX@{}}
\toprule
Component & Path \\
\midrule
C1/C2 edit engine, reorder null & \texttt{phaseN/c1\_edits.py} \\
\method probe (all arms) & \texttt{phaseN/probe\_cafe.py} \\
Paired bootstrap contrasts & \texttt{phaseN/paired\_contrasts.py} \\
Ordering, per-$\sigma$, variance decomposition & \texttt{phaseN/analyze\_cafe.py} \\
COCO reference pool & \texttt{phaseN/build\_coco\_pool.py} \\
Reference similarity, PCA, CKA, SVCCA, RSA & \texttt{phaseN/ref\_similarity.py} \\
CKA / debiased CKA / SVCCA with self-tests & \texttt{phaseN/repsim\_metrics.py} \\
Projector training & \texttt{phaseN/train\_repa\_projector.py} \\
\midrule
\method results & \texttt{phaseN/cafe\_probe\_105300/} \\
Reorder null & \texttt{phaseN/reorder\_null\_105323/} \\
Reference similarity (+\,\texttt{repa}) & \texttt{phaseN/refsim\_repa\_106517/} \\
Projector & \texttt{phaseN/repa\_projector\_106522/} \\
\bottomrule
\end{tabularx}
\end{table}

Running \texttt{python3 phaseN/repsim\_metrics.py} executes the CKA and SVCCA self-tests quoted in
\Cref{sec:repmetrics}.

\end{document}
